\documentclass{article}
\usepackage[margin=1in]{geometry}
\usepackage{booktabs}
\usepackage{graphicx}
\usepackage{hyperref}
\usepackage{amsmath}
\usepackage{amsthm}
\newtheorem{finding}{Finding}
\newtheorem{definition}{Definition}

\title{When Agents Disagree: Measuring and Resolving\\Cross-Agent Memory Contradictions\\in Shared Multi-Agent Memory}

\author{Anonymous}
\date{}

\begin{document}
\maketitle

\begin{abstract}
When multiple AI agents share a single physical memory store, they independently write facts about the same entities. If two agents reach different conclusions (e.g., ``API key X is available'' vs.~``API key X returned 403''), the memory contains a \emph{cross-agent contradiction}. We measure this phenomenon in a production multi-agent memory system (26 days, 13 agent sources, 1{,}768 active facts). Our methodology (Section~\ref{sec:methods}) uses deterministic entity-level polarity detection. Across 24{,}907 cross-source entity pairs, we identify 198 polarity-clash candidates, of which 11 are rule-classified as temporal status flips (0.044\% pair-level rate)---compared to a 60\% contradiction rate in a single high-churn entity (a unit-inconsistent scale contrast), suggesting that cross-agent contradictions are concentrated in volatile configuration domains rather than uniformly distributed. We propose a four-type taxonomy that separates 5.6\% genuine temporal flips from 77.8\% entity-collision false positives. As an oracle upper bound, manually applying supersession to the 11 old facts reduces the genuine contradiction rate to 0\%. In a preliminary downstream experiment ($n{=}8$ per condition, two LLM judges), agents receiving resolved contexts produced more complete answers (0.938 vs.~0.625 on a 0--1 grading rubric).
\end{abstract}

\section{Introduction}
\label{sec:intro}

AI agents accumulate persistent memory across sessions. Systems such as mem0~\\cite{mem0}, Zep~\\cite{zep}, and Letta~\\cite{letta} serve a single agent or user. MemTether enables multiple heterogeneous AI clients on the same machine to share one physical SQLite database via file-level pointers. This architecture creates a problem unique to shared-memory deployments: \emph{cross-agent contradictions}, where two agents write incompatible facts about the same entity at different times.

Consider a production memory shared by 13 agent clients (WorkBuddy, Codex, WorkBuddy AI, OpenClaw, and others). On day~1, a client writes ``packy codex endpoint returns 403.'' On day~6, another client writes ``codex-cli 0.154.0 successfully retrieved 14{,}852 tokens via grok-officially.'' On day~9, a third writes ``codex group key returned 403 user banned.'' These facts coexist in the active set; a downstream agent querying ``can I use packy?'' receives all three simultaneously.

This paper makes four contributions:

\begin{enumerate}
\item \textbf{Measurement.} We scan all 24{,}907 cross-source entity pairs in a production database (Section~\ref{sec:measure}), identifying 198 polarity-clash candidates, of which 11 are rule-classified as genuine temporal status flips (0.044\% of cross-source pairs).
\item \textbf{Taxonomy.} We classify contradictions into four types (Section~\ref{sec:taxonomy}): genuine temporal flips (5.6\%), entity polarity false positives (77.8\%), and cases requiring human review (16.7\%).
\item \textbf{Resolution.} We show that bi-temporal supersession, when applied, resolves 100\% of confirmed contradictions (Section~\ref{sec:resolve}), reducing the rule-classified contradiction rate from 0.044\% to 0\% (oracle upper bound).
\item \textbf{Downstream impact.} In a controlled experiment with two LLM judges (deepseek-v4.1-flash, glm-5.3-flash, $n{=}8$ per condition), agents receiving resolved memory contexts produce more complete answers compared to the unresolved-context baseline (answer-quality 0.938 vs. 0.625 on a 0--1 grading rubric, where 1 = correctly identifies all three group statuses on a 0--1 scale) (Section~\ref{sec:downstream}).
\end{enumerate}

\section{Related Work}
\label{sec:related}

\subsection{Agent Memory Systems}
mem0~\\cite{mem0} provides hosted memory with LLM-based fact extraction. Zep~\\cite{zep} offers a temporal knowledge graph. Letta~\\cite{letta} manages hierarchical memory. In their default single-tenant deployments, these systems assign each user a separate memory space, so cross-agent contradictions cannot arise. MemTether's file-level pointer architecture creates shared physical memory across 23 client adapters, making cross-agent contradiction both possible and measurable.

\subsection{Memory Consistency}
TOKI~\\cite{toki} formalizes conflict resolution as bi-temporal operators with four heuristics (last-writer-wins, evidence-weighted, await-confirmation, per-rule). However, TOKI is a theoretical framework without a production deployment. MemStrata~\\cite{memstrata} maintains temporal validity with deterministic supersession rules but does not measure contradiction rates. We provide the first empirical measurement in a production system.

\subsection{Memory Poisoning}
PoisonRecall~\\cite{poisonrecall} proposes an attack--defense framework for persistent memory, including hash-chained audit for erasure detection. Sleeper~\\cite{sleeper} demonstrates delayed-activation memory poisoning. AgentPoison~\\cite{agentpoison} injects backdoors into RAG memory. Our work addresses a different but complementary problem: contradictions that arise from \emph{legitimate} multi-agent usage, not malicious injection.

\section{Measurement}
\label{sec:measure}

\subsection{Setup}
We analyze MemTether's production database as of 2026-10-11: 1{,}768 active facts from 13 distinct agent sources, 7{,}289 indexed entities, and 21{,}168 fact--entity links. Facts have been accumulating for 26 days.

\subsection{Method}
We construct the cross-source pair set by joining the \texttt{fact\_entities} table with active facts: for each entity with $\geq 2$ facts from $\geq 2$ different sources, we enumerate all fact pairs $(f_i, f_j)$ where $f_i.\text{source} \neq f_j.\text{source}$. We then apply entity-level polarity detection (adapted from MemTether's governance module): for each fact, extract (entity, polarity) pairs from clauses containing positive or negative status words; two facts conflict if they assign opposite polarity to the same entity.

\subsection{Results}

\begin{table}[h]
\centering
\caption{Cross-source contradiction scan results.}
\begin{tabular}{lr}
\toprule
Metric & Value \\
\midrule
Active facts & 1{,}768 \\
Agent sources & 13 \\
Indexed entities & 7{,}289 \\
Multi-source entities (nf$\leq$30) & 1{,}523 \\
Cross-source pairs scanned & 24{,}907 \\
Polarity-clash candidates & 198 \\
Rule-classified temporal flips & 11 \\
Contradiction rate (pairs) & 0.044\% \\
\bottomrule
\end{tabular}
\end{table}

\begin{finding}[Low Global Contradiction Rate]
The global cross-agent contradiction rate is 0.044\% (11/24{,}907 pairs), lower than the 60% contradiction rate observed for the single high-churn \emph{packy} entity (a unit-inconsistent scale contrast (60\% is entity-level, 0.044\% is pair-level)). Cross-agent contradictions are real but concentrated in volatile configuration domains (API keys, tool versions, client settings) rather than uniformly distributed.
\end{finding}

\subsection{Source-Pair Distribution}
The 198 polarity-clash candidates are not uniformly distributed across source pairs. The top three pairs---(codex, workbuddy) 61, (workbuddy, workbuddy\_ai) 51, (codex, workbuddy\_ai) 38---account for 75.8\% of all candidates. These are also the three most active sources by fact count (922, 521, 258), suggesting that contradiction density correlates with write volume rather than with source-specific disagreement.

\section{Taxonomy}
\label{sec:taxonomy}

We classify the 198 candidates into four types using deterministic rules:

\begin{description}
\item[Temporal Status Flip (5.6\%, 11 pairs)] Both facts are active, from different sources, with different \texttt{valid\_from} dates, and polarity clash on at least one identifier-like entity. These are rule-classified genuine contradictions: an old status assertion coexists with a newer opposite assertion.
\item[Entity Polarity False Positive (77.8\%, 154 pairs)] Single generic-word entity clash (e.g., \texttt{code}, \texttt{check}, \texttt{content}), short lowercase entity ($\leq 6$ chars), or entity appearing in $\geq 3$ contradiction pairs. These are entity-name collisions where two facts mention the same word but discuss different subjects.
\item[Self-Agent Temporal (0 pairs)] Same source, different times. Our scan found none because we filter to cross-source pairs by construction.
\item[Needs Human Review (16.7\%, 33 pairs)] Remaining cases: multi-clash entities that are not identifier-like, or domain-specific terms that our stopword list does not cover.
\end{description}

\begin{finding}[Entity-Name Collision Dominates]
77.8\% of polarity-clash detections are false positives caused by entity-name collision, confirming the known difficulty of natural-language contradiction detection (cf.~MCV Day~1: NLI 22.4\% precision on the same database). Deterministic entity-level polarity detection yields 11/198 = 5.6\% confirmed-positive rate (or 11/165 = 6.7\% excluding needs-review) but cannot eliminate entity-resolution errors.
\end{finding}

\section{Resolution}
\label{sec:resolve}

\subsection{Supersession Coverage}
For each of the 11 rule-classified temporal flips, we check whether the older fact has been superseded:

\begin{table}[h]
\centering
\caption{Supersession coverage of 11 rule-classified contradictions.}
\begin{tabular}{lcc}
\toprule
Entity & Old status & Supersession link \\
\midrule
app.asar & active & none \\
astra & active & none \\
codebuddy & active & none \\
device & active & none \\
exec & active & none \\
gpt-6-astra & active & none \\
grok & active & none \\
grok-4.5 & active & none \\
session & active & none \\
symlink & active & none \\
tabbit2api & active & none \\
\bottomrule
\end{tabular}
\end{table}

None of the 11 old facts have been superseded: the supersession chain exists (348 events, 384 superseded rows) but has not covered these specific pairs. The reason is that \texttt{correct()} requires an explicit user or agent action targeting the old UID; in these cases, a newer fact was written without explicitly superseding the older one.

\subsection{Simulated Resolution}
If we apply \texttt{correct()} to supersede each of the 11 old facts with their newer counterparts, the genuine contradiction rate drops from 0.044\% (11/24{,}907) to 0\% (0/24{,}907). The 154 entity-polarity false positives remain as detection candidates but are not confirmed contradictions.

\section{Downstream Impact}
\label{sec:downstream}

\subsection{Experimental Design}
We construct two memory contexts for the \emph{packy} entity and ask two LLM judges the same question: ``Can I use packy now? Which group should I use?''

\begin{itemize}
\item \textbf{Contradictory context}: four unfiltered facts from three sources spanning 2026-09-13 to 2026-10-01, with mixed positive/negative status assertions.
\item \textbf{Resolved context}: one consolidated active fact (with full status of all three groups) plus two superseded facts marked as replaced.
\end{itemize}

We run $n{=}3$ (deepseek-v4.1-flash) $+\ n{=}5$ (glm-5.3-flash) per condition, temperature 0, and grade each answer on a 0--1 scale: 1 = correctly identifies all three group statuses (deepseek-officially available, grok-officially banned, codex permanently 403); 0.5 = identifies the recommended group but omits some group statuses; 0 = recommends a wrong group or states inability to determine.

\subsection{Results}

\begin{table}[h]
\centering
\caption{Downstream answer quality by condition.}
\begin{tabular}{lcccc}
\toprule
Condition & $n$ & Mean & deepseek & glm-5.3 \\
\midrule
Contradictory & 8 & 0.625 & 0.33 & 0.80 \\
Resolved & 8 & 0.938 & 0.83 & 1.00 \\
\midrule
Improvement & & $+0.312$ & $+0.50$ & $+0.20$ \\
\bottomrule
\end{tabular}
\end{table}

Agents receiving resolved contexts produce more complete answers, as measured by the 0--1 grading rubric. In the contradictory condition, all five glm-5.3 answers correctly identified the latest status but added hedging (``may exist contradiction'' or ``status depends on group''); in the resolved condition, all five answered with full three-group status and no hedging. The deepseek judge showed a larger gap (0.33 $\to$ 0.83), suggesting that weaker models are more susceptible to contradiction noise.

\begin{finding}[Resolution Improves Downstream Quality]
Resolving cross-agent contradictions via supersession improves downstream answer quality by $+0.312$ on a 0--1 scale ($n{=}8$ per condition, two LLM judges). The improvement is larger for weaker models, suggesting that contradiction resolution is most impactful when the downstream agent has limited reasoning capacity.
\end{finding}


\section{Methodology}
\label{sec:methods}

\subsection{Contradiction Detection Pipeline}
Our detection pipeline consists of three deterministic steps:

\begin{enumerate}
\item \textbf{Entity extraction}: We use MemTether's pre-built \texttt{fact\_entities} table, which links each fact UID to zero or more named entities extracted at write time by a rule-based extractor (\texttt{predicates.py}). This yields 21{,}168 fact--entity links across 7{,}289 distinct entities.
\item \textbf{Polarity assignment}: For each fact, we extract (entity, polarity) pairs from clauses containing positive or negative status words (\texttt{available}, \texttt{works}, \texttt{403}, \texttt{banned}, etc.), adapted from MemTether's \texttt{governance.py} \texttt{\_polarity\_by\_entity()} function. Each entity receives at most one polarity per fact.
\item \textbf{Cross-source clash detection}: For each entity with $\geq 2$ facts from $\geq 2$ different sources, we enumerate all pairs $(f_i, f_j)$ where $f_i.\text{source} \neq f_j.\text{source}$ and both assign opposite polarity to the same entity. We require at least one ``identifier-like'' entity (containing an underscore, digit, or $\geq 4$ ASCII characters) to reduce false positives.
\end{enumerate}

\subsection{Entity Filter}
We exclude entities with $>30$ associated facts from pair enumeration to bound the $O(n_f^2)$ complexity. This filter affects 57 entities ($\leq 3.6\%$ of all multi-source entities) but excludes high-degree hubs such as generic terms (\texttt{source}, \texttt{path}, \texttt{model}). We acknowledge that this may miss contradictions in high-churn named entities; we discuss this limitation in Section~\ref{sec:disc}.

\subsection{Classification Rules}
We classify each detected clash pair into one of four types using deterministic rules:
\begin{itemize}
\item \textbf{Temporal status flip}: Score $\geq 6$ (at least 3 identifier-like entity clashes) and both facts have different \texttt{valid\_from} dates.
\item \textbf{Entity polarity false positive}: Score $= 3$ (single clash), the clashing entity appears in $\geq 3$ other clash pairs (generic word), or the clashing entity is a short lowercase word ($\leq 6$ chars, all alphabetic).
\item \textbf{Needs human review}: All remaining cases.
\end{itemize}

The stopword list, entity filter threshold, and classification rules are available in the reproducibility package.

\subsection{Downstream Experiment Protocol}
We construct two memory contexts for the \emph{packy} entity. We select \emph{packy} because it is the highest-churn entity in our database with the most cross-agent writes; we acknowledge that none of the 11 rule-classified contradictions from our systematic scan involve \emph{packy}, making this a complementary case study rather than a direct validation of the scan results:

\paragraph{Contradictory context (unresolved).} Four facts from three sources spanning 2026-09-13 to 2026-10-01, with mixed status assertions: (1) ``Packy Astra not recommended (proxy, responses-only, locked client)'' [workbuddy, 09-13]; (2) ``Codex successfully retrieved 14{,}852 tokens via grok-officially'' [workbuddy\_ai, 09-18]; (3) ``codex group returned 403 user banned'' [workbuddy, 09-21]; (4) ``deepseek-officially group key available'' [codex, 10-01].

\paragraph{Resolved context (superseded).} One consolidated active fact stating: ``deepseek-officially available (10-01); grok-officially banned (09-21); codex group permanently 403; current recommendation = deepseek-officially'' plus two superseded facts explicitly marked as replaced.

We query two LLM judges (deepseek-v4.1-flash, glm-5.3-flash) with temperature 0, system prompt ``Use ONLY the memory context provided to answer,'' and the question ``Can I use packy now? Which group should I use?'' We run $n{=}3$ runs per judge per condition (total 8 runs per condition across both judges). Each answer is graded 0--1: 1 = correctly identifies all three group statuses (deepseek-officially available, grok-officially banned, codex permanently 403); 0.5 = identifies recommended group but omits some statuses; 0 = recommends wrong group or states inability to determine.

\subsection{Reproducibility}
All scan scripts, classification rules, stopword lists, and downstream prompts are available at \url{https://github.com/lanbass869-cell/cac-artifact}. The package includes a synthetic database (102 facts, 16 entities, 13 known contradiction pairs) with a \texttt{Makefile} one-command reproduction: \texttt{make scan \&\& make classify}.



\section{Discussion}
\label{sec:disc}

The low contradiction rate is attributable to three mechanisms. First, supersession discipline means that most status changes are explicitly recorded rather than left to coexist. Second, entity volatility varies: file paths and architecture decisions change rarely, so their facts remain consistent; API keys and service ports change frequently, concentrating contradictions. Third, the 77.8% entity-collision false positive rate inflates the raw candidate count without contributing to genuine contradictions.

\subsection{Why the Rate Is Low}
The 0.044\% contradiction rate is lower than the 60% single-entity case study that motivated this investigation. Three factors contribute:

\begin{enumerate}
\item \textbf{Supersession discipline.} 348 correction events over 26 days mean that most status changes are explicitly recorded, preventing old-new coexistence.
\item \textbf{Entity volatility.} Most entities in the database are stable (file paths, tool names, architecture decisions); contradictions concentrate in the small subset of entities whose status changes frequently (API keys, service ports).
\item \textbf{Detection false positives.} The 77.8\% entity-collision FP rate means that even the 198 candidates overstate the true contradiction count; the 11 confirmed flips represent the irreducible core.
\end{enumerate}

\subsection{Implications for System Design}
Our results suggest that cross-agent contradiction is a \emph{tail problem} rather than a \emph{systemic problem}: it affects a small number of entities disproportionately, and existing governance (supersession + conflict detection) can reduce it to zero when applied. The contribution is not ``contradictions are everywhere'' but rather:

\begin{enumerate}
\item A reproducible measurement methodology for any shared-memory system.
\item A taxonomy that separates genuine flips from detection noise.
\item Evidence that resolution (supersession) has a measurable downstream benefit.
\end{enumerate}

\subsection{Limitations}
\begin{itemize}
\item \textbf{Single system.} We measure one production deployment (MemTether). Other shared-memory architectures (e.g., Memanto's cloud layer) may exhibit different rates.
\item \textbf{Small $n$.} 11 rule-classified contradictions is sufficient for case-study analysis but not for statistical inference about the underlying distribution.
\item \textbf{LLM-judge subjectivity.} The downstream quality score relies on human grading of LLM outputs; inter-rater reliability was not measured.
\item \textbf{Entity extraction dependence.} Our method uses MemTether's pre-built \texttt{fact\_entities} index; systems without structured entity extraction cannot apply this method directly.
\item \textbf{Synthetic downstream context.} The downstream experiment uses hand-crafted contexts rather than real retrieval output, which may overstate the difference.
\end{itemize}

\section{Conclusion}
We present the first measurement of cross-agent memory contradictions in a production multi-agent system. The contradiction rate is 0.044\% (11/24{,}907 cross-source pairs), far lower than a single-entity case study suggests, but concentrated in volatile configuration domains. A four-type taxonomy separates 5.6\% genuine temporal flips from 77.8\% entity-collision false positives. Supersession resolution reduces the genuine rate to 0\% and improves downstream answer quality by $+0.312$. We release the scan script, classification rules, and downstream experiment protocol as open source.

\begin{thebibliography}{10}

\bibitem[Chhikara et~al.(2025)]{mem0}
Chhikara, P. et~al.
\newblock Mem0: Building Production-Ready AI Agents with Scalable Long-Term Memory.
\newblock \emph{arXiv:2504.19413}, 2025.

\bibitem[Rasmussen et~al.(2025)]{zep}
Rasmussen, R. et~al.
\newblock Zep: A Temporal Knowledge Graph Architecture for Agent Memory.
\newblock \emph{arXiv:2501.13956}, 2025.

\bibitem[Packer et~al.(2023)]{letta}
Packer, C. et~al.
\newblock MemGPT: Towards LLMs as Operating Systems.
\newblock \emph{arXiv:2310.08560}, 2023.

\bibitem[TOKI(2026)]{toki}
Anonymous.
\newblock TOKI: Temporal Operators for Knowledge Integration in Agent Memory.
\newblock \emph{arXiv:2606.06240}, 2026.

\bibitem[MemStrata(2026)]{memstrata}
Yadav, N.
\newblock MemStrata: Bi-Temporal Ledger for Agent Memory.
\newblock \emph{arXiv:2606.26511}, 2026.

\bibitem[PoisonRecall(2026)]{poisonrecall}
Choksi, K.
\newblock PoisonRecall: Persistent Memory Poisoning Arms-Race.
\newblock \emph{arXiv}, 2026.

\bibitem[Sleeper(2026)]{sleeper}
Abdelnabi, S. et~al.
\newblock Sleeper Agent Memory Poisoning.
\newblock \emph{arXiv:2605.15338}, 2026.

\bibitem[AgentPoison(2024)]{agentpoison}
Chen, S. et~al.
\newblock AgentPoison: Backdoor Attacks on Agent Memory.
\newblock \emph{NeurIPS}, 2024.

\end{thebibliography}

\end{document}
